% ===== 原文 pp.42–46（PDF p.52–56） =====

% [原文 p.42]
\chapter{Vafa-Witten 理论与 GL-扭转理论在圆上的紧化}

我们开始执行把 $\cN=4$、$d=4$ 的 SYM 理论的各个扭转版本紧化到更低维 TFT 的计划，先从在圆上紧化到三维的情形入手。我们的目标有两重：我们希望借此阐明低维 TFT（特别是，4d TFT 之间的对偶关系会延拓为它们的维数约化版本之间的对偶关系）；此外，低维 TFT 也将帮助我们理解原来 4d TFT 的若干侧面。

例如，按照第 4 章讨论的维数约化技巧，在圆上紧化所得的 3d 理论的边界条件，将分类原来 4d 理论中面算子的谱（我们把对这些边界条件的分析推迟到下一章）。

在本章中，我们把 Vafa-Witten 理论以及 $t=\ii$、$t=1$、$t=0$ 的 GL-扭转理论在圆上约化到三维。所得 3d TFT 将被证明是拓扑的、带规范场的 $\sigma$-模型：$t=\ii$ 处的约化给出 [14] 中引入的 Rozansky-Witten 模型的带规范场版本，而 $t=1$ 与 $t=0$ 处的约化则给出 [6] 中引入的 3d A-模型 TFT 的带规范场版本；Vafa-Witten 理论在圆上紧化时与 $t=0$ 理论等价。

% [原文 p.43]
\section{拓扑规范理论的紧化}

在认真着手 $S^1$ 约化之前，让我们先对拓扑规范理论的紧化作几点一般性评论。考虑我们的某个 4d TFT 定义在乘积流形 $M=W\times S^1$ 上，其中 $W$ 是黎曼的、有向的三维流形，$S^1$ 是周长为 $2\pi R$ 的圆。我们希望找到 $R\to0$ 极限下的有效描述。（更一般地，可以用维数不低于二的紧致紧化流形来紧化。）

场论紧化的通常 Kaluza–Klein 概念是：把 4d 理论的场沿圆作傅里叶模展开，并在 $R\to0$ 极限下丢弃除常数模之外的所有模，理由是非零模的质量会趋于无穷。以此方式，得到 $W$ 上的有效 3d 理论。（更确切地说，必须把非零模从路径积分中积掉，这可能给有效理论带来量子修正。）

% [Tutor Note: p.43 的 Kaluza–Klein 论证可直接核对：沿 $S^1$ 的第 $n$ 个傅里叶模的 4d 质量满足 $m_n^2\sim n^2/R^2$，故 $R\to0$ 时 $n\ne0$ 的模质量发散而退耦。本章取 $2\pi R=2\pi$ 即 $R=1$，此时常数模质量为零、$n\ne0$ 模质量 $\sim n$；括号中的“量子修正”正来自把这些模从路径积分积掉所生的行列式贡献。]
% [Tutor Note: p.43 的规范变换取 $g(x^4)=\exp(\ii mx^4)$；因 $x^4$ 的周期为 $2\pi$，该变换单值当且仅当 $m\in\Z$。按本章的数学约定（联络反 Hermite），$A_4\mapsto A_4+g^{-1}\partial_4g=A_4+\ii m$；于是 $\exp(-2\pi A_4)$ 在整数平移下不变，这正是应当用 $h$ 而非 $A_4$ 来写 3d 理论的根源。]

但 TFT“不在意”$R$ 的大小，意思是：$R$ 的改变只以 BRST-正合的量改变作用量，因而使拓扑关联函数保持不变。因此（模去我们稍后讨论的微妙之处）约化后的理论——尽管它描述的是场空间中的一个简化区域，其中场不沿紧致方向变化——与原来理论严格等价。这一便利事实已被用于 [1]：通过在两维曲面上紧化，$t=\ii$ 与 $t=1$ 的 GL-扭转理论获得了 2d 描述。在本章余下部分，我们把周长取为方便的值 $2\pi R=2\pi$。

当我们的理论是拓扑规范理论时，有两个重要的微妙之处使这一图景复杂化。第一个是：乘积流形上的给定场位形可能在紧化流形上留下一些未被破缺的残余连续规范对称性。这类场位形对应于有效 $\sigma$-模型描述靶空间中的奇点：在这些点上，紧化理论中会打开定性上新的分支。实际上，人们能做的只是把这些位形“手动”切除（例如施加适当的边界条件以排除它们的出现）。[1] 中的 2d 有效理论正是以这种方式被截断的。在目前的情形下这个问题不会出现，因为我们的紧化流形只是 $S^1$。

然而第二个微妙之处对我们的分析很重要。先考虑规范场为 $A$ 的阿贝尔规范理论。直接宣布 $A$ 与 $x^4$ 无关、并令分量 $A_4$ 为实标量场，并不完全正确，其中 $x^4$ 是 $S^1$ 的坐标，从 $x^4=0$ 到 $x^4=2\pi$。这是因为 $x^4$-依赖的规范变换 $g(x^4)=\exp(\ii mx^4)$ 会把 $A_4$ 的值平移一个整数：
\[
    A_4 \mapsto A_4 + \ii m, \qquad m \in \Z .
\]
因此，如果直接用标量 $A_4$ 来写，3d 有效理论是定义不良的。

问题的解决办法是改用周期标量
\[
    h = \exp(-2\pi A_4)
\]
来写有效理论；该场在 $x^4$-依赖的规范变换下不变。

% [原文 p.44]
对非阿贝尔规范群，还有一个额外的复杂性：要求 4d 理论的场不依赖 $x^4$ 坐标并不是规范不变的条件。一种做法是先固定一个规范使 $A_4$ 不依赖 $x^4$；确实，这在 $A_4=0$ 附近，即 $A$ 沿 $S^1$ 的和乐接近单位元时是可行的。但一般而言，“$A_4$ 与 $x^4$ 无关”这一条件并不能消去作 $x^4$-依赖规范变换的自由度。例如，设 $A_4$ 正比于 $\fg$ 中满足
\[
    \exp(2\pi\mu) = 1
\]
的元素 $\mu$。这样的 $\mu$ 恰恰是那些落在 $G$ 的余特征标格（cocharacter lattice）的 $G$-轨道中的元素。于是规范变换
\[
    g(x^4) = \exp(\mu x^4)
\]
把 $A_4$ 平移 $\mu$：
\[
    A_4 \mapsto A_4 + \mu .
\]
这样的规范变换一般而言会使其他场带上 $x^4$-依赖性。

% [Tutor Note: p.44 的 $\exp(2\pi\mu)=1$ 说明 $\mu$ 落在极大环面的整格（余特征标格）中，从而 $g(x^4)=\ee^{\mu x^4}$ 单值。对 $G=U(1)$ 即 $\mu=\ii m$，此时 $A_4\mapsto A_4+\ii m$ 恰好是第 43 页的整数平移，而和乐 $h=\ee^{-2\pi A_4}$ 不变——这是检验“大规范变换”作用的最快锚点。]
% [Tutor Note: p.44 的 $h\in\Omega^0(W,P_W\times_GG)$ 来自 $h\mapsto g^{-1}hg$：规范变换把绕圈的平行移动整体共轭，故 $h$ 是与 $P_W$ 相伴、纤维为 $G$、转移函数取共轭作用的丛的截面，而非 $G$ 值函数。这正是非阿贝尔情形不能把 $h$ 当普通标量场的根源，也是后文 $\Tr\{h^{-1}D_rh\,h^{-1}D^rh\}$ 中 $h^{-1}(\cdot)h$ 结构的原因。]

解决办法仍然是：在理论的描述中用群取值的场 $h$ 取代 $A_4$，$h$ 现在定义为 $A$ 绕紧致方向的和乐。让我们看看这在细节上如何运作。在通过紧化计算 4d 理论的某个关联函数时，方便的做法是先把所有插入算子的支集形变到嵌入超曲面 $W_0\subset M$ 上，$W_0$ 定义为 $W\times S^1$ 中 $x^4=0$ 的点构成的轨迹。（由于拓扑不变性，这种形变不影响关联函数的值。）接着，我们定义 3d 场 $h$ 为联络绕从点 $(\vec x,0)\in W_0$ 出发、绕紧致方向一周的圈的和乐：
\[
    h(\vec x) = P\exp\left(-\int_0^{2\pi}\dd x^4 A_4(\vec x,x^4)\right)
\]
其中 $\vec x\in W$；和乐在 $N$ 维定义表示中计算。设 $P_W$ 为 $P$ 沿嵌入映射
\[
    W \longrightarrow W_0 \longrightarrow M
\]
拉回得到的 $W$ 上的主 $G$-丛。

对非阿贝尔的 $G$，和乐 $h$ 不再是规范不变的；确切地说，它是规范协变的，作为与 $P_W$ 相伴、典型纤维为 $G$、由 $G$ 在自身上的共轭作用给出的纤维丛的截面而变换；也就是说，
\[
    h \in \Omega^0(W, P_W \times_G G) .
\]
特别地，$h$ 是三维中一个完全良定义的标量场。

% [原文 p.45]
\section{Vafa-Witten 理论的 \texorpdfstring{$S^1$}{S1} 紧化}

先考虑把 Vafa-Witten 理论放在乘积流形 $M=W\times S^1$ 上。我们将证明，在对 $M$ 的这一特定几何上，该理论与靶空间为 $X=\fg$ 的 3d A-模型的带规范场版本等价，其中规范扇区是对 B.2.1 节讨论的 3d A-型规范理论的一个轻微修改（把李代数取值的标量场 $\varsigma$ 换成群取值的场）。

\para{约化后的场含量}

我们以一种能显示各场在 3d 约化理论中将扮演何种角色的方式来重新标记 Vafa-Witten 理论的场。以下指标 $r,s,t=1,2,3$ 指 $W$ 上的坐标，上标 $[VW]$ 表示 Vafa-Witten 理论中相应的场，如我们在第 3 章中所写。场含量的 3d 约化将分为规范扇区与物质扇区。3d 规范扇区将来自重新标记后的场：
\begin{gather*}
    A_r = A_r^{[VW]}, \qquad \lambda_r = \psi_r^{[VW]} \\
    \varsigma = A_4^{[VW]}, \qquad \rho = \psi_4^{[VW]} \\
    \sigma = \sigma^{[VW]}, \\
    \bar\sigma = \bar\sigma^{[VW]}, \qquad \tilde\rho = \eta^{[VW]} \\
    \tilde\lambda_r = \chi_{r4}^{[VW]}, \qquad H_r = H_{r4}^{[VW]} .
\end{gather*}
3d 物质扇区将来自重新标记后的场：
\begin{gather*}
    \phi = \phi^{[VW]}, \qquad \eta = \zeta^{[VW]} \\
    b_r = B_{r4}^{[VW]}, \qquad \psi_r = \tilde\psi_{r4}^{[VW]} \\
    \beta = \tilde\chi_4^{[VW]}, \qquad P = \tilde H_4^{[VW]} \\
    \chi_r = \tilde\chi_r^{[VW]}, \qquad \tilde P_r = \tilde H_r^{[VW]} .
\end{gather*}

% [原文 p.46]
Vafa-Witten 理论的 BRST 变分 (3.8) 用重新标记后的场写出后，成为
\begin{gather}
    \delta_QA = \lambda, \qquad \delta_Q\lambda = -d_A\sigma, \notag\\
    \delta\varsigma = \rho, \qquad \delta_Q\rho = [\sigma,\varsigma] \notag\\
    \delta_Q\sigma = 0, \notag\\
    \delta_Q\bar\sigma = \tilde\rho, \qquad \delta_Q\tilde\rho = [\sigma,\bar\sigma] \notag\\
    \delta_Q\tilde\lambda = H, \qquad \delta_QH = [\sigma,\tilde\lambda] \notag\\
    \delta_Q\phi = \eta, \qquad \delta_Q\eta = [\sigma,\phi] \notag\\
    \delta_Qb = \psi, \qquad \delta_Q\psi = [\sigma,b] \notag\\
    \delta_Q\beta = P, \qquad \delta_QP = [\sigma,\beta] \notag\\
    \delta_Q\chi = \tilde P, \qquad \delta_Q\tilde P = [\sigma,\chi] . \tag{5.1}
\end{gather}

% [Tutor Note: p.46 式 (5.1) 第三行原文写作 $\delta\varsigma=\rho$，缺下标 $Q$；由其余各行及第 3 章式 (3.8) 的 BRST 变分可判定应为 $\delta_Q\varsigma=\rho$。译文照原文保留 $\delta$，勿据此作别的解释。]
% [Tutor Note: p.46 的 $\delta_Q^2h=[\sigma,h]$ 可由同页两式直接验证：$h^{-1}\delta_Qh=2\pi\hat\rho$ 给出 $\delta_Qh=2\pi h\hat\rho$，而 $\delta_Q(h\hat\rho)=(2\pi)^{-1}[\sigma,h]$，故 $\delta_Q^2h=2\pi\,\delta_Q(h\hat\rho)=[\sigma,h]$，恰为参数 $\sigma$ 的规范变换；其中 $2\pi$ 因子恰好与和乐归一化（周长 $2\pi$）相消。]
% [Tutor Note: p.46 的第一行 $\delta_QA=\lambda$、$\delta_Q\lambda=-d_A\sigma$ 中 $A,\lambda$ 未标指标，指 3d 1-形式整体；与第 3 章式 (3.8) 对照时注意那里 $A$ 为 4d 联络、此处已按 $r,s,t=1,2,3$ 截断。]

如果我们天真地只保留上述场的常数傅里叶模以约化到 3d，那么我们将得到靶空间为 $\fg$（由场 $\phi$ 参数化）的带规范场的 3d A-模型 TFT 的正确场含量与变分，如 B.2 节所写。然而，如上文所讨论，场 $A_4$ 作为 3d 标量场是定义不良的，必须在理论的描述中用它的和乐 $h$ 取代。此外，我们还必须把它的 BRST 伙伴 $\rho$ 换成一个适当平均的场。我们定义
\[
    \hat\rho = -\frac{1}{2\pi}\int_0^{2\pi}\dd x^4\, H^{-1}\rho H
\]
其中
\[
    H(x^4) = P\exp\left(-\int_0^{x^4}\dd t\, A_4\right)
\]
是 $A$ 沿从 $0$ 到 $x^4$ 的圆弧的和乐，从而特别地 $H(2\pi)=h$。于是人们可以验证（需要一些计算），$h$ 与 $\hat\rho$ 的 BRST 变分由
\[
    h^{-1}\delta_Qh = 2\pi\hat\rho, \qquad
    \delta_Q(h\hat\rho) = (2\pi)^{-1}[\sigma,h] .
\]
给出。结合这两式可以确认 $\delta_Q^2h=[\sigma,h]$，即以 $\sigma$ 为参数的规范变换。场 $\hat\rho$ 按截面变换：
\[
    \hat\rho \in \Omega^0(W, \operatorname{ad} P_W) .
\]
